\section{Introduction}
\label{sec:intro}

Spatial transcriptomics preserves the tissue coordinates of gene-expression measurements, allowing direct study of microenvironments, intercellular signalling, and spatially patterned disease processes that dissociated single-cell RNA sequencing (scRNA-seq) cannot access. Since its recognition as \emph{Nature Methods}' Method of the Year in 2020 \citep{marx2021method}, the technology has become routine in both basic and translational research. The most widely adopted platform, 10x Genomics Visium, operates at a resolution of 55~\textmu m per spot, with each spot typically containing 1--10 cells. This ``spot mixing'' problem means that the measured expression at each location is a superposition of signals from multiple cells of potentially different types, obscuring the cell-type-specific biology that researchers seek to characterise.

The gap between spot-level measurements and single-cell resolution has driven the development of computational methods that recover single-cell information from mixed observations. These methods vary in ambition: some estimate cell-type proportions within each spot, others place individual reference cells at spatial coordinates, and the most aggressive synthesise complete per-cell expression profiles at sub-spot resolution. They differ in their assumptions, data requirements, and the confidence warranted by their outputs. Existing reviews cover parts of this space. \citet{gaspardboulinc2025deconvolution} provide a thorough survey of cell-type deconvolution but do not extend to mapping or reconstruction. Earlier reviews by \citet{longo2021integrating} and \citet{tian2023expanding} predate pseudo single-cell reconstruction methods entirely. No existing review treats the complete computational pipeline from spots to cells as a unified progression.

This review provides the first systematic treatment of the full four-stage pipeline for enhancing spatial transcriptomics resolution. The four stages are spot-level deconvolution, which estimates cell-type proportions per spot; single-cell mapping, which assigns reference cells to spatial coordinates; pseudo single-cell reconstruction, which synthesises per-cell expression at sub-spot resolution; and multi-modal enhancement, which uses histology and paired modalities to refine earlier outputs. We survey more than 25 methods published between 2020 and 2026, propose a two-axis taxonomy based on output resolution and information source, and provide a practical decision tree for method selection. Critically, we delineate what pseudo single-cell reconstruction can and cannot reliably infer, addressing the growing concern that reconstructed cells are being treated as observational data in downstream analyses without appropriate caveats.

The remainder of this review is organised as follows. Section~\ref{sec:background} introduces spatial transcriptomics platforms and formalises the resolution gap. Sections~\ref{sec:stage1}--\ref{sec:stage4} present the four stages in detail. Section~\ref{sec:comparison} provides unified comparison and practical guidance. Section~\ref{sec:limitations} discusses open challenges, and Section~\ref{sec:future} outlines future directions.
